\documentclass[10pt,a4paper]{amsart}
\usepackage[margin=19mm]{geometry}
\usepackage{amsmath,amssymb,mathtools}
\usepackage[scr=rsfs,cal=euler]{mathalfa}
\usepackage{xfrac}
\usepackage[unicode,hidelinks,bookmarks=false]{hyperref}
\newcommand{\ie}{i.e.\ }
\newcommand{\cf}{cf.\ }
\newcommand{\resp}{respectively\ }
\newcommand{\Subsection}{\subsection}
\providecommand{\nobreakdash}{\nobreak\hbox{-}}
\setlength{\emergencystretch}{2em}
\multlinegap0pt
\usepackage{cancel}
\usepackage{mathtools}
\usepackage{algorithm}\usepackage{algpseudocode}
\usepackage{bm}
\usepackage{tikz}
\usepackage{amssymb}
\newtheorem{assump}{Assumption}
\newtheorem{lem}{Lemma}
\newtheorem{prop}{Proposition}
\newcommand\B{\mathcal B} %OK
	\newcommand\BD{\end{color}}
\newcommand\DB{\begin{color}{black}}
\DeclareMathOperator*{\Grad}{\boldsymbol\nabla}
\renewcommand\H{\mathbf{H}}
\newcommand\HuOm{\mathbf{H}^1(\Omega)} %OK
\newcommand\Hub{\mathbf{H}^1(\B)} %OK
\newcommand\LL{{\boldsymbol{\Lambda}}} %OK
\newcommand\Ldo{L^2_0(\Omega)} %OK
\newcommand\Pcal{\mathcal{P}}
\renewcommand\S{\mathbf{S}} %OK
\newcommand\Sline{\S} %OK
\newcommand\T{\mathcal{T}} %OK
\newcommand\Vline{\mathbf{V}} %OK
\newcommand\X{\mathbf{X}} %OK
\newcommand\Xbar{\overline{\X}} %OK
\newcommand{\abs}[1]{\left\lvert#1\right\rvert}
\newcommand\ds{\mathrm{d}\s} %OK
\DeclareMathOperator{\grads}{\boldsymbol\nabla_s}
\newcommand\mmu{{\boldsymbol\mu}} %OK
\newcommand{\norm}[1]{\left\lVert#1\right\rVert}
\newcommand\quadnode{\mathbf{p}}
\newcommand\quadweigth{\omega}
\newcommand\reftri{\widehat{\tri}} %OK
\newcommand\s{\mathbf{s}} %OK
\newcommand\tri{\mathrm{T}} %OK
\newcommand\trifj{\tri_{\mathrm{f},j}} %OK
\renewcommand\v{\mathbf{v}} %OK
\newcommand\w{\mathbf{w}} %OK
\title{Quadrature error estimates on non--matching grids\\in a fictitious domain framework\\for fluid--structure interaction problems}
\author{Daniele Boffi, Fabio Credali, Lucia Gastaldi}
\date{Source: arXiv:2406.03981v2 (CC BY 4.0)}
\begin{document}
\maketitle

\noindent\emph{Source and licence.} Quadrature error estimates on non--matching grids\\in a fictitious domain framework\\for fluid--structure interaction problems. Version arXiv:2406.03981v2 (CC BY 4.0), distributed by arXiv under the Creative Commons Attribution 4.0 International licence, http://creativecommons.org/licenses/by/4.0/, as recorded in the arXiv licence metadata for this exact version and shown on the abstract page https://arxiv.org/abs/2406.03981v2. Full licence notice: \texttt{licenses/arxiv-2406.03981-CC-BY-4.0.txt}.\par\smallskip

\noindent\textit{Editorial scope.} Counted unit: the article's prop prop:l2grad\_estimate (source0.tex lines 1256--1271). The article's complete original proof of it is delivered with it (source0.tex lines 1273--1377, one \texttt{proof} environment of 105 lines). No mathematics is written, repaired or re-derived: the statement and the proof are the article's own lines, in the article's own notation, and no retained line is substituted or re-flowed.  Retained uncounted as the source-local context the proof needs: lines 629--637; lines 1025--1029; lines 1033--1040; lines 1080--1095; lines 1103--1110; lines 1165--1180.  Nothing was omitted from any retained span, so the package carries no omission record.  No retained line contains a citation, so the wrapper carries no bibliography and no bibliography entry.  No retained line contains a figure environment, an image-inclusion command or drawing code, so no image is delivered and the package declares no asset dependency.  Editorial formatting only, in addition: an AMS \texttt{amsart} preamble that restates the article's own declarations and macros used by the retained lines, namely the environments \texttt{assump}, \texttt{lem}, \texttt{prop} on separate counters, together with the standard packages those lines need.  The retained environments print in this document as Assump 1, Lemma 1, Proposition 1 in order of appearance, whereas the article prints the same items with the numbering of its own full document.  The complete native source of the article as distributed in the arXiv e-print for this exact version is bundled unchanged as \texttt{sources/160383/source0.tex}.
	\begin{equation}\label{eq:element}
		\begin{aligned}
			\Vline_h &= \{ \v\in\HuOm: \v_{\mid
				\tri}\in[\Pcal_1(\tri)]^2 \quad \forall \tri\in\T_{h/2}^\Omega \} \\
			Q_h &= \{q\in \Ldo\cap H^1(\Omega): q_{\mid \tri}\in \Pcal_1(\tri)\quad\forall \tri\in\T_{h}^\Omega\}\\
			\S_h &= \{\w\in\Hub:\w_{\mid \tri}\in [\Pcal_1(\tri)]^2\quad\forall \tri\in\mathcal{T}_{h}^\B\}\\
			\LL_h&=\{\mmu\in\Hub:\mmu_{\mid \tri}\in [\Pcal_1(\tri)]^2\quad\forall \tri\in\mathcal{T}_{h}^\B\}.
		\end{aligned}
	\end{equation}

	\begin{assump}\label{assum:xbar}	
		We assume that $\Xbar\in\Sline_h$. This implies that it is linear in $\tri$ for all $\tri\in\T_h^B$. Consequently, the composed function
		$\v_h(\Xbar)$ is continuous and piecewise linear on
		$\tri$.
	\end{assump} 

	\begin{lem}
		Let $\tri\in\T_h^\B$ be such that $\Xbar(\tri)$ is not included in an element of
		$\T_{h/2}^\Omega$ and let Assumption~\ref{assum:xbar} hold true. Then $\v_h(\Xbar)\in \H^{1+s}(\tri)$ for $0\le s<1/2$ and 
		\[
		\|\grads\v_h(\Xbar)\|_{s,\tri}\le \frac{C}{1-2s} \|\grads\v_h(\Xbar)\|_{0,\tri}.
		\]
		\label{le:duran}
	\end{lem}

	\begin{prop}\label{prop:l2_quad_error}
		Let $\Vline_h$ and $\LL_h$ be given by \eqref{eq:element}. Let us assume that $\Xbar$ satisfies Assumption~\ref{assum:xbar}.
		Given a quadrature rule $\{(\quadnode_k^0,\quadweigth_k^0)\}_{k=1}^{K_0}$ exact for quadratic polynomials, i.e. such that
		\begin{equation}\label{eq:quadr}
			\widehat{E}(\widehat{f}) = 0 \quad\forall\widehat{f}\in\Pcal_2(\reftri),
		\end{equation}
		the following estimate holds true %for $s\in[0,1/2)$
		\begin{equation}\label{eq:l2_quad_error}
			\abs{E_\B(\mmu_h\cdot\v_h(\Xbar))}
			\le
			C h_\B^{3/2}|\log h_\B^{min}|
			\norm{\mmu_h}_{0,\B}\DB\abs{\v_h}_{1,\Omega}\BD
			\quad\forall\mmu_h\in\LL_h,\forall\v_h\in\Vline_h,
		\end{equation}
		where $E_\B$ is the sum of $E_\tri(\mmu_h\cdot\v_h(\Xbar))$ for all $\tri\in\T_h^\B$ and  $h_\B^{min}=\min_{\tri\in\T_h^\B} h_\tri$. 
	\end{prop}

		\begin{equation}\label{eq:solid_mesh_partition}
			\begin{aligned}
				&\T_{h,1}^\B = \{\tri\in\T_h^\B:\DB\Xbar(\tri)\BD\text{ is included in an element of }
				\T_{h/2}^\Omega\}
				\\
				&\T_{h,2}^\B = \T_h^\B \setminus \T_{h,1}^\B.
			\end{aligned}
		\end{equation}

		\begin{equation}\label{eq:3rd_part1_l2}
			\begin{aligned}
				&\abs{\abs{\tri} \sum_{k=1}^{K_0}\quadweigth_k^0 \mmu_h(\quadnode_k^0)\cdot
					\big(\v_h(\Xbar(\quadnode_k^0))-\v_I(\quadnode_k^0)\big)}\\
				&\qquad \le \bigg(\abs{\tri}
				\sum_{k=1}^{K_0}\quadweigth_k^0\,\mmu_h(\quadnode_k^0)^2 \bigg)^{1/2} \,
				\bigg(\abs{\tri}
				\sum_{k=1}^{K_0}\quadweigth_k^0\,
				\big(\v_h(\Xbar(\quadnode_k^0))-\v_I(\quadnode_k^0)\big)^2\bigg)^{1/2}\\
				&\qquad =\bigg(\int_\tri \abs{\mmu_h}^2\bigg)^{1/2}\, \bigg(\abs{\tri}
				\sum_{k=1}^{K_0}\quadweigth_k^0\,
				\big(\v_h(\Xbar(\quadnode_k^0))-\v_I(\quadnode_k^0)\big)^2\bigg)^{1/2}\\
				&\qquad \le K_0\abs{\tri}^{1/2}\norm{\mmu_h}_{0,\tri}
				\norm{\v_h(\Xbar)-\v_I}_{\infty,\tri}
			\end{aligned}
		\end{equation}

	\begin{prop}\label{prop:l2grad_estimate}
		Let $\Vline_h$ and $\LL_h$ given by \eqref{eq:element}. Let us assume that $\Xbar$ satisfies Assumption~\ref{assum:xbar}.
		Given a quadrature rule $\{(\quadnode_k^1,\quadweigth_k^1)\}_{k=1}^{K_1}$ exact for constants, i.e. such that
		\begin{equation}
			\widehat{E}(\widehat{f}) = 0
			\quad\forall\widehat{f}\in\Pcal_0(\reftri),
		\end{equation}
		and a  quasi-uniform mesh  $\T_h^\Omega$, the following estimate holds
		\begin{equation}\label{eq:h1_quad_error}
			\abs{E_\B(\grads\mmu_h:\grads\v_h(\Xbar))}
			\le
			C\bigg( h_\B^{1/2}|\log h_\B^{min}| + \frac{h_\B}{h_\Omega} \bigg)
			\norm{\grads \mmu_h}_{0,\B}\norm{\Grad\v_h}_{0,\Omega}
		\end{equation}
		for all $\mmu_h\in\LL_h$ and $\v_h\in\Vline_h$. Here $E_\B$ is the sum of $E_\tri(\grads\mmu_h:\grads\v_h(\Xbar))$ for all $\tri\in\T_h^\B$ and  $h_\B^{min}=\min_{\tri\in\T_h^\B} h_\tri$. 
	\end{prop}

	\begin{proof}
		Using the same technique as in the proof of Proposition~\ref{prop:l2_quad_error},
		we partition the solid mesh $\T_h^\B$ as in
		\eqref{eq:solid_mesh_partition} and we work locally in an element
		$\tri\in\T_{h,2}^\B$: this means that we have to bound the difference
		\begin{equation}
			E_\tri(\grads\mmu_h:\grads\v_h(\Xbar)) = \int_\tri
			\grads\mmu_h:\grads\v_h(\Xbar)\,\ds -
			\abs{\tri}\sum_{k=1}^{K_1}\quadweigth_k^1\grads\mmu_h(\quadnode_k^1):
			\grads\v_h(\Xbar(\quadnode_k^1)).
		\end{equation}
		We recall that $\grads\mmu_h\in[\Pcal_0(\tri)]^2$, while $\grads\v_h(\Xbar)$ is
		a discontinuous piecewise constant function in~$\tri$.
		Working again with the interpolant $\v_I\in[\Pcal_1(\tri))]^2$ of $\v_h(\Xbar)$,
		we can write
		\begin{equation}\label{eq:manipulation_h1}
			\begin{aligned}
				E_\tri(\grads\mmu_h:\grads\v_h(\Xbar)) = &\int_\tri \grads\mmu_h:\big(\grads\v_h(\Xbar)-\grads\v_I\big)\,\ds\\
				&+ \int_\tri \grads\mmu_h:\grads\v_I\,\ds - \abs{\tri}\sum_{k=1}^{K_1}\quadweigth_k^1\grads\mmu_h(\quadnode_k^1):\grads\v_I(\quadnode_k^1)\\
				&\DB-\BD  \abs{\tri} \sum_{k=1}^{K_1}\quadweigth_k^1 \grads\mmu_h(\quadnode_k^1): \big(\grads\v_h(\Xbar(\quadnode_k^1))-\grads\v_I(\quadnode_k^1)\big)
			\end{aligned}
		\end{equation}
		so that we can deal separately with each term.
		
		For the first term, we use the Cauchy--Schwarz inequality
		\begin{equation}
			\abs{\int_\tri
				\grads\mmu_h:\big(\grads\v_h(\Xbar)-\grads\v_I\big)\,\ds}\le
			\norm{\grads\mmu_h}_{0,\tri} \norm{\grads\v_h(\Xbar)-\grads\v_I}_{0,\tri},
		\end{equation}
		then using Lemma~\ref{le:duran} we get
		\[
		\norm{\grads\v_h(\Xbar)-\grads\v_I}_{0,\tri}
		\le h_\tri^s\norm{\grads\v_h(\Xbar)}_{s,\tri}
		\le \frac{h_\tri^s}{1-2s}\norm{\grads\v_h(\Xbar)}_{0,\tri}
		\]
		so that, taking $s=\frac12+\frac1{\log h_\tri}$, we find
		\begin{equation}
			\begin{aligned}
				&\abs{\int_\tri \grads\mmu_h:\big(\grads\v_h(\Xbar)-\grads\v_I\big)\,\ds}\\
				&\hspace{3cm}\le Ch_\tri^{1/2}\abs{\log h_\tri} \norm{\grads\mmu_h}_{0,\tri}\norm{\grads\v_h(\Xbar)}_{0,\tri}.
			\end{aligned}
		\end{equation}
		
		The second term is zero by construction of $\v_I$ and the choice of quadrature rule
		\begin{equation}
			E_\tri(\grads\mmu_h:\grads\v_I) = \int_\tri \grads\mmu_h:\grads\v_I\,\ds - \abs{\tri}\sum_{k=1}^{K_1}\quadweigth_k^1\grads\mmu_h(\quadnode_k^1):\grads\v_I(\quadnode_k^1)=0.
		\end{equation}
		
		For the last term, we proceed as in \eqref{eq:3rd_part1_l2} so that we have
		\begin{equation}\label{eq:3rd_part1_h1}
			\begin{aligned}
				&\abs{\abs{\tri} \sum_{k=1}^{K_1}\quadweigth_k^1 \grads\mmu_h(\quadnode_k^1): \big(\grads\v_h(\Xbar(\quadnode_k^1))-\grads\v_I(\quadnode_k^1)\big)}\\
				&\hspace{3.5cm}\le K_1 \abs{\tri}^{1/2} \norm{\grads\mmu_h}_{0,\tri}\norm{\grads\v_h(\Xbar)-\grads\v_I}_{\infty,\tri}%\\
			\end{aligned}
		\end{equation}
		
		In order to estimate $\norm{\grads\v_h(\Xbar)-\grads\v_I}_{\infty,\tri}$,
		we take into account the discontinuity of $\grads\v_h(\Xbar)$ in $\tri$:
		hence we consider again the tessellation of $\tri$ into disjoint polygons
		$P_j$, $j=1,\dots,J$, with $\tri=\bigcup_{j=1}^J P_j$
		and $\grads\v_h(\Xbar)$ constant in each $P_j$.
		Using that $\norm{\grads\v_I}_{\infty,\tri} \le\norm{\grads\v_h(\Xbar)}_{\infty,\tri} $ ,
		we can write
		\begin{equation}\label{eq:step_polygons}
			\begin{aligned}
				\norm{\grads\v_h(\Xbar)-\grads\v_I}_{\infty,\tri} 
				&= \max_{j=1,\dots,J} \norm{\grads\v_h(\Xbar)-\grads\v_I}_{\infty,P_j}\\
				&\le 2 \max_{j=1,\dots,J} \norm{\grads\v_h(\Xbar)}_{\infty,P_j}.
			\end{aligned}
		\end{equation}
		We notice that each polygon $P_j$ is the inverse image
		through $\Xbar$ of $\Xbar(\tri)\cap\trifj$ with $\trifj\in\T_{h/2}^\Omega$.
		Moreover, we denote by $\omega_\tri$ the macroelement
		consisting of all fluid triangles intersecting $\Xbar(\tri)$,
		i.e. ${\omega_\tri = \bigcup_{j=1}^J \trifj}$.
		Therefore, by applying the inverse inequality
		$\|\w_h\|_{\infty,\tau}\le C_I/h_\tau\|\w_h\|_{0,\tau}$ for $\w_h\in\Vline_h$
		and $\tau\in\T_{h/2}^\Omega$, we get
		%
		\begin{equation}\label{eq:step_macroelement}
			\begin{aligned}
				\max_{j=1,\dots,J} \norm{\grads\v_h(\Xbar)}_{\infty,P_j}
				&\le \DB C_{\Xbar}\BD\max_{j=1,\dots,J} \norm{\Grad\v_h}_{\infty,\trifj}\\
				&\le \DB C_{\Xbar}\BD\max_{j=1,\dots,J}\,\frac{\DB C_{\trifj}\BD}{h_{\trifj}}
				\norm{\Grad\v_h}_{0,\trifj}
				\le \frac{\DB C\BD}{h_{\Omega}} \norm{\Grad\v_h}_{0,\omega_\tri},
			\end{aligned}
		\end{equation}
		%
		\DB where $C_{\Xbar}$ depends on $\|\Xbar\|_{W^{1,\infty}(\B)}$ and\BD, in the last
		inequality, we exploited that $\T_{h/2}^\Omega$ is quasi-uniform. Putting
		together \eqref{eq:3rd_part1_h1} with \eqref{eq:step_polygons} and
		\eqref{eq:step_macroelement}, we end up with
		\begin{equation}
			\begin{aligned}
				&\abs{\abs{\tri} \sum_{k=1}^{K_1}\quadweigth_k^1 \grads\mmu_h(\quadnode_k^1): \big(\grads\v_h(\Xbar(\quadnode_k^1))-\grads\v_I(\quadnode_k^1)\big)}\\
				& \hspace{1.5cm} \le C \frac{\abs{\tri}^{1/2}}{h_{\Omega}} \norm{\grads\mmu_h}_{0,\tri} \norm{\Grad\v_h}_{0,\omega_\tri}\
				\le C\frac{h_\B}{h_{\Omega}} \norm{\grads\mmu_h}_{0,\tri} \norm{\Grad\v_h}_{0,\omega_\tri}.
			\end{aligned}
		\end{equation}
		%
		Finally, summing on over all solid elements and exploiting $\Xbar(\B)\subset\Omega$, the global estimate \eqref{eq:h1_quad_error} is proven.
		
	\end{proof}

\end{document}
